% GENERATED FILE. Edit canonical lessons, then run npm run generate:books.
% canonical-source-hash: fnv1a64:3075cb666b7e6561
% canonical-lessons: GE-C332-begruenden, GE-C332-begruendung, GE-C332-naemlich, GE-C332-zumal, GE-C332-erlaeutern

\chapter{Gründe nennen}
\label{ch:ge-grunde-nennen}
\hlchaptermodality{332}

\begin{chapteropening}
\textbf{By the end of this chapter:} \emph{I can give reasons and explain them: justify a choice, add the real reason with nämlich and an extra one with zumal.}

Complete the last lesson of chapter 332: I can give reasons and explain them: justify a choice, add the real reason with nämlich and an extra one with zumal.
\end{chapteropening}

\section[begründen]{begründen — to give reasons for, to justify; to found}

\label{lesson:GE-C332-begruenden}



\begin{quote}
Before the new one: say the German for to talk, then the German for to argue.
\end{quote}

\subsection*{You'll want to know: begründen}
\textbf{begründen} — \textquotedblleft{}to give reasons for, to justify; to found\textquotedblright{}.

\begin{quote}
Der Lehrer hat seine Entscheidung gut begründet.
\end{quote}

\begin{quote}
\emph{The teacher gave good reasons for his decision.}
\end{quote}

\emph{Etwas begründen} — to back it up with reasons. It also means to found: \emph{eine Schule begründen}.

\begin{cousinweb}
\emph{begründen} is inherited from Middle High German \emph{begründen} \textquotedblleft{}to found, to make solid\textquotedblright{} (EtymDB), built on \emph{Grund} \textquotedblleft{}ground, reason\textquotedblright{}, from Proto-Germanic \emph{grunduz} \textquotedblleft{}ground\textquotedblright{}: you put ground under your words.
\end{cousinweb}

\begin{grammarlens}[title={What you've built}]
One more word for telling, explaining and describing at B1.
\end{grammarlens}

\subsection*{Your turn}
\begin{itemize}
\raggedright
  \item \emph{Say it:} \emph{begründen}
  \item \emph{Say it:} \emph{begründen}, once more
  \item \emph{Say it:} a sentence of your own with \emph{begründen}
  \item \emph{Recall:} say the German for to talk, then the German for to argue, then say \emph{begründen} again
\end{itemize}

\begin{tcolorbox}[breakable,colback=teal!4,colframe=teal!35!black,title={Before you move on}]
What does begründen mean? (\textquotedblleft{}To give reasons for, to justify; to found\textquotedblright{}.) Say it once more.
\end{tcolorbox}
\section[die Begründung]{die Begründung — the justification, the reasons given}

\label{lesson:GE-C332-begruendung}



\begin{quote}
Before the new one: say the German for to argue, then the German for to give reasons for.
\end{quote}

\subsection*{You'll want to know: die Begründung}
\textbf{die Begründung} — \textquotedblleft{}the justification, the reasons given\textquotedblright{}.

\begin{quote}
Ohne Begründung hat er die Einladung abgelehnt.
\end{quote}

\begin{quote}
\emph{He turned down the invitation without giving any reason.}
\end{quote}

Plural: \emph{die Begründungen}. \emph{Mit der Begründung, dass …} — on the grounds that …

\begin{cousinweb}
\emph{Begründung} is the noun of \emph{begründen}, formed with \emph{-ung}. The verb is built on \emph{Grund} \textquotedblleft{}ground\textquotedblright{}, inherited from Proto-Germanic \emph{grunduz} (EtymDB).
\end{cousinweb}

\begin{grammarlens}[title={What you've built}]
One more word for telling, explaining and describing at B1.
\end{grammarlens}

\subsection*{Your turn}
\begin{itemize}
\raggedright
  \item \emph{Say it:} \emph{die Begründung}
  \item \emph{Say it:} \emph{die Begründung}, once more
  \item \emph{Say it:} a sentence of your own with \emph{die Begründung}
  \item \emph{Recall:} say the German for to argue, then the German for to give reasons for, then say \emph{die Begründung} again
\end{itemize}

\begin{tcolorbox}[breakable,colback=teal!4,colframe=teal!35!black,title={Before you move on}]
What does die Begründung mean? (\textquotedblleft{}The justification, the reasons given\textquotedblright{}.) Say it once more.
\end{tcolorbox}
\section[nämlich]{nämlich — you see (adding the reason); namely}

\label{lesson:GE-C332-naemlich}



\begin{quote}
Before the new one: say the German for to give reasons for, then the German for the justification.
\end{quote}

\subsection*{You'll want to know: nämlich}
\textbf{nämlich} — \textquotedblleft{}you see (adding the reason); namely\textquotedblright{}.

\begin{quote}
Ich komme heute nicht. Ich bin nämlich krank.
\end{quote}

\begin{quote}
\emph{I am not coming today. I am ill, you see.}
\end{quote}

\emph{Nämlich} adds the reason after the fact and stands after the verb, not at the start: \emph{Ich bin nämlich krank.} It also means namely: \emph{zwei Farben, nämlich Rot und Blau}.

\begin{cousinweb}
\emph{nämlich} goes back to Middle High German \emph{nemelīch} \textquotedblleft{}by name, expressly\textquotedblright{}, built on \emph{Name} \textquotedblleft{}name\textquotedblright{}. EtymDB has no reliable entry for the adverb, so this follows the usual account: first you named the thing, later the reason.
\end{cousinweb}

\begin{grammarlens}[title={What you've built}]
One more word for telling, explaining and describing at B1.
\end{grammarlens}

\subsection*{Your turn}
\begin{itemize}
\raggedright
  \item \emph{Say it:} \emph{nämlich}
  \item \emph{Say it:} \emph{nämlich}, once more
  \item \emph{Say it:} a sentence of your own with \emph{nämlich}
  \item \emph{Recall:} say the German for to give reasons for, then the German for the justification, then say \emph{nämlich} again
\end{itemize}

\begin{tcolorbox}[breakable,colback=teal!4,colframe=teal!35!black,title={Before you move on}]
What does nämlich mean? (\textquotedblleft{}You see (adding the reason); namely\textquotedblright{}.) Say it once more.
\end{tcolorbox}
\section[zumal]{zumal — especially since, all the more as}

\label{lesson:GE-C332-zumal}



\begin{quote}
Before the new one: say the German for the justification, then the German for you see.
\end{quote}

\subsection*{You'll want to know: zumal}
\textbf{zumal} — \textquotedblleft{}especially since, all the more as\textquotedblright{}.

\begin{quote}
Wir bleiben zu Hause, zumal es draußen regnet.
\end{quote}

\begin{quote}
\emph{We are staying at home, especially since it is raining outside.}
\end{quote}

\emph{Zumal} starts a clause with the verb at the end, like \emph{weil}. It adds a second, extra reason.

\begin{cousinweb}
\emph{zumal} was built from \emph{zu} \textquotedblleft{}at\textquotedblright{} and \emph{Mal} \textquotedblleft{}time, occasion\textquotedblright{}, from Old High German \emph{māl} (EtymDB), and is usually said to have first meant \textquotedblleft{}at the same time, all together\textquotedblright{}, from which the sense \textquotedblleft{}especially since\textquotedblright{} grew; EtymDB has no entry for the compound.
\end{cousinweb}

\begin{grammarlens}[title={What you've built}]
One more word for telling, explaining and describing at B1.
\end{grammarlens}

\subsection*{Your turn}
\begin{itemize}
\raggedright
  \item \emph{Say it:} \emph{zumal}
  \item \emph{Say it:} \emph{zumal}, once more
  \item \emph{Say it:} a sentence of your own with \emph{zumal}
  \item \emph{Recall:} say the German for the justification, then the German for you see, then say \emph{zumal} again
\end{itemize}

\begin{tcolorbox}[breakable,colback=teal!4,colframe=teal!35!black,title={Before you move on}]
What does zumal mean? (\textquotedblleft{}Especially since, all the more as\textquotedblright{}.) Say it once more.
\end{tcolorbox}
\section[erläutern]{erläutern — to explain in detail, to clarify}

\label{lesson:GE-C332-erlaeutern}



\begin{quote}
Before the new one: say the German for you see, then the German for especially since.
\end{quote}

\subsection*{You'll want to know: erläutern}
\textbf{erläutern} — \textquotedblleft{}to explain in detail, to clarify\textquotedblright{}.

\begin{quote}
Die Ärztin hat mir die Ergebnisse genau erläutert.
\end{quote}

\begin{quote}
\emph{The doctor explained the results to me in detail.}
\end{quote}

More formal than \emph{erklären}. Inseparable: \emph{sie erläutert, sie hat erläutert}.

\begin{cousinweb}
\emph{erläutern} is built on \emph{lauter} \textquotedblleft{}clear, pure\textquotedblright{}, inherited from Middle High German \emph{lūter}, Old High German \emph{hlūttar} and Proto-Germanic \emph{hlūtraz} \textquotedblleft{}clear, clean\textquotedblright{} (EtymDB): you make a thing clear.
\end{cousinweb}

\begin{grammarlens}[title={What you've built}]
One more word for telling, explaining and describing at B1.
\end{grammarlens}

\subsection*{Your turn}
\begin{itemize}
\raggedright
  \item \emph{Say it:} \emph{erläutern}
  \item \emph{Say it:} \emph{erläutern}, once more
  \item \emph{Say it:} a sentence of your own with \emph{erläutern}
  \item \emph{Recall:} say the German for you see, then the German for especially since, then say \emph{erläutern} again
\end{itemize}

\begin{tcolorbox}[breakable,colback=teal!4,colframe=teal!35!black,title={Before you move on}]
What does erläutern mean? (\textquotedblleft{}To explain in detail, to clarify\textquotedblright{}.) Say it once more.
\end{tcolorbox}
